\tag{LR2}
\]
Only this \emph{sign} cochain is asserted to take root-of-unity values;
the full scalar splitting also involves the $a_\sigma$.

On $G_L$, $b$ is a conjugation-invariant character. Outside $T$ its
inertia characters are quadratic. Only finitely many are nontrivial;
twist by the quadratic character of the product of those odd primes
to kill them. At each odd prime there is a unique nontrivial quadratic
inertia character. This changes no square or coboundary in (LR2),
and introduces possible ramification only at $2\in T$.
The kernel field $L'$ of $b|G_L$ is normal over $\mathbb Q$ and
\[
 [L':\mathbb Q]\le2m[L:\mathbb Q]\le2304R.
\tag{LR3}
\]

\paragraph{The compatible modular representations.}
Choose a number field containing all normalization scalars, the values
of $b$, and a splitting field of $B$. For each coefficient prime $\ell$
and an embedding into $\overline{\mathbb Q}_\ell$, set
\[
 \rho_\ell(\sigma)=b(\sigma)^{-1}a_\sigma\mu_\sigma\circ\sigma.
\]
The cocycle identity makes this a representation commuting with $B$;
after splitting $B$, matrix units give $\rho_\ell=r_\ell\oplus r_\ell$
with $\dim r_\ell=2$. On $G_{L'}$, $\rho_\ell$ is the Tate action of
$A$. Faltings' semisimplicity and endomorphism theorems and $\operatorname{End}^0(A)=B$
show that the multiplicity representation is absolutely irreducible.
Proper-smooth comparison and finite descent, as in
Lemma~\ref{lem:de-rham-descent}, make $r_\ell$ de Rham with Tate
Hodge--Tate weights $\{0,-1\}$; Berger's monodromy theorem makes it
potentially semistable.

It is odd. On Betti homology complex conjugation followed by the
holomorphic isogeny $\mu_c$ exchanges the two Hodge summands, each
a simple $B\otimes\mathbb C$-module. Its action on the multiplicity
space is therefore nonscalar. After the scalar correction its square
is one, so its eigenvalues are $1,-1$. The representation is unramified
outside $T\cup\{\ell\}$, by (LR2) and the reduction and field properties.
